% Options for packages loaded elsewhere
\PassOptionsToPackage{unicode}{hyperref}
\PassOptionsToPackage{hyphens}{url}
\documentclass[
]{article}
\usepackage{xcolor}
\usepackage{amsmath,amssymb}
\setcounter{secnumdepth}{-\maxdimen} % remove section numbering
\usepackage{iftex}
\ifPDFTeX
  \usepackage[T1]{fontenc}
  \usepackage[utf8]{inputenc}
  \usepackage{textcomp} % provide euro and other symbols
\else % if luatex or xetex
  \usepackage{unicode-math} % this also loads fontspec
  \defaultfontfeatures{Scale=MatchLowercase}
  \defaultfontfeatures[\rmfamily]{Ligatures=TeX,Scale=1}
\fi
\usepackage{lmodern}
\ifPDFTeX\else
  % xetex/luatex font selection
\fi
% Use upquote if available, for straight quotes in verbatim environments
\IfFileExists{upquote.sty}{\usepackage{upquote}}{}
\IfFileExists{microtype.sty}{% use microtype if available
  \usepackage[]{microtype}
  \UseMicrotypeSet[protrusion]{basicmath} % disable protrusion for tt fonts
}{}
\makeatletter
\@ifundefined{KOMAClassName}{% if non-KOMA class
  \IfFileExists{parskip.sty}{%
    \usepackage{parskip}
  }{% else
    \setlength{\parindent}{0pt}
    \setlength{\parskip}{6pt plus 2pt minus 1pt}}
}{% if KOMA class
  \KOMAoptions{parskip=half}}
\makeatother
\usepackage{longtable,booktabs,array}
\newcounter{none} % for unnumbered tables
\usepackage{calc} % for calculating minipage widths
% Correct order of tables after \paragraph or \subparagraph
\usepackage{etoolbox}
\makeatletter
\patchcmd\longtable{\par}{\if@noskipsec\mbox{}\fi\par}{}{}
\makeatother
% Allow footnotes in longtable head/foot
\IfFileExists{footnotehyper.sty}{\usepackage{footnotehyper}}{\usepackage{footnote}}
\makesavenoteenv{longtable}
\setlength{\emergencystretch}{3em} % prevent overfull lines
\providecommand{\tightlist}{%
  \setlength{\itemsep}{0pt}\setlength{\parskip}{0pt}}
\usepackage{bookmark}
\IfFileExists{xurl.sty}{\usepackage{xurl}}{} % add URL line breaks if available
\urlstyle{same}
\hypersetup{
  pdftitle={DPM Gauge Embedding: SO(2)\_DPM = Light-Cone Plane in SO(26) (G3 Closure)},
  pdfauthor={Daniel T. Murphy},
  hidelinks,
  pdfcreator={LaTeX via pandoc}}

\title{DPM Gauge Embedding: SO(2)\_DPM = Light-Cone Plane in SO(26) (G3
Closure)}
\author{Daniel T. Murphy}
\date{2026-05-10}

\begin{document}
\maketitle

\section{\texorpdfstring{PAPER\_1163 -- DPM Gauge Embedding:
\(SO(2)_{\rm DPM}\) = Light-Cone Plane in
\(SO(26)\)}{PAPER\_1163 -- DPM Gauge Embedding: SO(2)\_\{\textbackslash rm DPM\} = Light-Cone Plane in SO(26)}}\label{paper_1163-dpm-gauge-embedding-so2_rm-dpm-light-cone-plane-in-so26}

\textbf{Author:} Daniel T. Murphy (daniel.murphy00@gmail.com)
\textbf{Framework:} UQFF v5.78 -- Star-Magic Physics \textbf{Source:}
Direct identification with PAPER\_050 light-cone gauge, standard bosonic
string textbook result. \textbf{Integration Date:} May 10, 2026 (Session
250) \textbf{Classification:} Lagrangian Gap Closure -- G3 of
\texttt{\_lagrangian\_rederivation\_outline.py}

\begin{center}\rule{0.5\linewidth}{0.5pt}\end{center}

\[\boxed{\;SO(26) \;\supset\; SO(24)_{\rm trans} \times SO(2)_{\rm DPM}\;, \qquad \dim(325) = \dim(276) + \dim(1) + 48_{\rm coset}\;}\]

The DPM gauge group \(SO(2)\) -- the CW/CCW counter-rotation that
generates \(F_{\rm DPM} = I\cdot A\cdot (\omega_1 - \omega_2)\) --
embeds in the full 26D Lorentz group \(SO(26)\) as the
\textbf{light-cone gauge 2-plane} already removed in textbook bosonic
string quantization. Zero free parameters: the embedding is the same one
used in
\href{whitepapers/PAPER_050_26D_Manifold_Compactification_3plus1_Spacetime.md\#L220}{PAPER\_050
§3}.

\begin{center}\rule{0.5\linewidth}{0.5pt}\end{center}

\subsection{Abstract}\label{abstract}

We close gap \textbf{G3} of the Lagrangian re-derivation outline by
identifying the SO(2) gauge structure of the DPM (di-pseudo-monopole
CW/CCW pair) with the \textbf{light-cone gauge 2-plane} that is removed
in standard bosonic string light-cone quantization at
\(D_{\rm crit}=26\). The DPM force
\(F_{\rm DPM} = I\cdot A\cdot(\omega_1-\omega_2)\) is literally an SO(2)
angular-momentum generator acting on the 2-dimensional fibre orthogonal
to the 24 transverse Lorentz dimensions. The embedding
\(SO(26)\supset SO(24)_{\rm trans}\times SO(2)_{\rm DPM}\) is
dimensionally exact (\(325 = 276 + 1 + 48_{\rm coset}\)) and Dynkin
index 1, requiring zero new free parameters beyond \(D_{\rm crit}=26\).

\begin{center}\rule{0.5\linewidth}{0.5pt}\end{center}

\subsection{1. The Gap}\label{the-gap}

The Lagrangian outline (\texttt{\_lagrangian\_rederivation\_outline.py},
L56-58) records:

\begin{quote}
\textbf{G3.} \(\mathcal{L}_{\rm DPM}\) uses an SO(2) gauge structure
(CW-CCW pair) but its embedding in the full 26D internal symmetry group
is not written down.
\end{quote}

The required statement is: identify a canonical embedding
\(\iota: SO(2) \hookrightarrow SO(26)\) with \textbf{zero free
parameters}, verify that the DPM current
\(F_{\rm DPM} = I\cdot A\cdot(\omega_1-\omega_2)\) is genuinely the
generator of this \(SO(2)\), and confirm the dimensional/index counting.

\subsection{2. The DPM Current Is an SO(2) Generator (codebase
canonical)}\label{the-dpm-current-is-an-so2-generator-codebase-canonical}

From
\href{COMPLETE_UQFF_EQUATIONS_REFERENCE.md\#L549}{COMPLETE\_UQFF\_EQUATIONS\_REFERENCE.md
L549} and
\href{source4_wolfram_COMPLETE_SPEC.md\#L664}{source4\_wolfram\_COMPLETE\_SPEC.md
L664}:

\[F_{\rm DPM} \;=\; I \cdot A \cdot (\omega_1 - \omega_2)\]

with \(\omega_1, \omega_2\) the angular velocities of the \textbf{CW}
and \textbf{CCW} rotating monopoles. This is exactly the form of an
angular momentum operator on a 2-plane:

\[T_{[12]} \;=\; -i(x_1 \partial_2 - x_2 \partial_1)\]

acting on the differential current \(I\cdot A\) with eigenvalue
\((\omega_1 - \omega_2)\). The CW/CCW pair are the two eigenstates of
\(T_{[12]}\) with eigenvalues \(\pm i\) -- the standard fundamental
representation of \(SO(2)\cong U(1)\).

Thus
\textbf{\(\mathcal{L}_{\rm DPM} = -\tfrac{1}{4}F_{[12]}^{\rm DPM} F^{[12]\rm DPM}\)
is a Yang-Mills Lagrangian for the gauge group \(SO(2)\) generated by
\(T_{[12]}\)}, with
\(F_{[12]}^{\rm DPM} = \partial_\mu A_\nu^{\rm DPM} - \partial_\nu A_\mu^{\rm DPM}\).
No free parameters.

\subsection{\texorpdfstring{3. The Canonical Embedding
\(SO(2) \hookrightarrow SO(26)\)}{3. The Canonical Embedding SO(2) \textbackslash hookrightarrow SO(26)}}\label{the-canonical-embedding-so2-hookrightarrow-so26}

In 26D Wick-rotated space (Euclidean signature for the closed-bosonic
worldsheet target), the full internal symmetry is \(SO(26)\).
\textbf{Textbook bosonic string light-cone quantization} (GSW Vol. 1 Ch.
2, Polchinski Vol. 1 §1.3) fixes a 2-plane gauge:

\[SO(26) \;\longrightarrow\; SO(24)_{\rm transverse} \,\times\, SO(2)_{\rm light\text{-}cone}\]

The 24 transverse dimensions carry the physical oscillator modes
\(\alpha_n^i\) with \(i=1,\dots,24\) (the ``magical \(24\)'' of the
bosonic string spectrum). The 2 light-cone dimensions \((X^+, X^-)\) are
gauge artefacts -- they parameterise the SO(2) plane removed by the
gauge fix.

PAPER\_050 explicitly states (§3): \textgreater{} \emph{``2 light-cone
gauge degrees removed from 26 = 24.''}

\textbf{G3 closure claim:} The SO(2) gauge group of the DPM is
\textbf{identically} this light-cone SO(2). The CW/CCW rotation
\(\omega_1 - \omega_2\) is the relative rotation of the two light-cone
coordinates \(X^\pm\) when viewed as a 2-plane in the original 26D
Lorentz frame.

\subsection{4. Dimensional Verification}\label{dimensional-verification}

The dimension counting must close exactly:

{\def\LTcaptype{none} % do not increment counter
\begin{longtable}[]{@{}
  >{\raggedright\arraybackslash}p{(\linewidth - 4\tabcolsep) * \real{0.3000}}
  >{\raggedleft\arraybackslash}p{(\linewidth - 4\tabcolsep) * \real{0.4000}}
  >{\raggedright\arraybackslash}p{(\linewidth - 4\tabcolsep) * \real{0.3000}}@{}}
\toprule\noalign{}
\begin{minipage}[b]{\linewidth}\raggedright
Group
\end{minipage} & \begin{minipage}[b]{\linewidth}\raggedleft
\(\dim G = n(n-1)/2\)
\end{minipage} & \begin{minipage}[b]{\linewidth}\raggedright
Role
\end{minipage} \\
\midrule\noalign{}
\endhead
\bottomrule\noalign{}
\endlastfoot
\(SO(26)\) & 325 & Full 26D Lorentz/internal \\
\(SO(24)_{\rm trans}\) & 276 & Transverse physical \\
\(SO(2)_{\rm DPM}\) & 1 & DPM gauge / light-cone \\
\textbf{Coset} \(SO(26)/(SO(24)\times SO(2))\) &
\textbf{\(325 - 276 - 1 = 48\)} & Off-diagonal generators \\
\end{longtable}
}

The coset dimension \(48 = 24 \times 2\) is the number of generators
that mix transverse \(i\) with light-cone \(\pm\) -- exactly the 24
\textbf{transverse \(\times\)} 2 \textbf{light-cone} off-diagonal
Lorentz boosts. This is textbook (cf.~GSW eq. 2.3.45). The number
factorises \textbf{identically} through the 24 light-cone physical modes
that carry the bosonic string spectrum.

\subsection{5. Dynkin Index of the
Embedding}\label{dynkin-index-of-the-embedding}

For an embedding \(\iota: H \to G\), the Dynkin index \(T(\iota)\)
measures the trace ratio of generators:
\[T(\iota) \;=\; \frac{\mathrm{Tr}_{R}\left(T^{H}_a T^{H}_b\right)}{\mathrm{Tr}_{R}\left(T^{G}_a T^{G}_b\right)}\]
For SO(2) embedded in SO(26) via the natural 2-plane (i.e., \(SO(2)\)
acts on the first two coordinates only), in the \textbf{vector
representation \(\mathbf{26}\) of SO(26)}:

\[T(SO(2) \to SO(26)) \;=\; 1\]

This is the \textbf{minimal possible index} (unit index), confirming
that the embedding is irreducible and does not pick up any ``level''
factor that would constitute a free parameter. Zero new parameters
introduced.

\subsection{\texorpdfstring{6. The Branching Rule for the Vector
\(\mathbf{26}\)}{6. The Branching Rule for the Vector \textbackslash mathbf\{26\}}}\label{the-branching-rule-for-the-vector-mathbf26}

Under \(SO(26) \supset SO(24)\times SO(2)\), the fundamental vector
representation \(\mathbf{26}\) branches as:
\[\mathbf{26} \;\longrightarrow\; (\mathbf{24},\,\mathbf{1}_0) \;\oplus\; (\mathbf{1},\,\mathbf{1}_{+1}) \;\oplus\; (\mathbf{1},\,\mathbf{1}_{-1})\]

where the subscript on the \(SO(2)\) singlet is the \textbf{U(1) charge}
\(q = \pm 1\). The two SO(2)-charged singlets \(\mathbf{1}_{\pm1}\) are
precisely the \textbf{CW (\(q=+1\)) and CCW (\(q=-1\)) monopoles} of the
DPM pair. The 24 SO(2)-neutral states are the transverse oscillator
modes that carry no DPM charge.

This is the \textbf{group-theoretic origin} of: * the CW/CCW pair (two
SO(2) charged states \(q=\pm 1\)); * the ``20-component'' +
``2-component'' DPM splitting referenced elsewhere in the framework; *
the \(\omega_1 - \omega_2\) structure as the SO(2) charge difference.

\subsection{7. Consistency With G6, G7,
G8}\label{consistency-with-g6-g7-g8}

{\def\LTcaptype{none} % do not increment counter
\begin{longtable}[]{@{}
  >{\raggedright\arraybackslash}p{(\linewidth - 4\tabcolsep) * \real{0.3333}}
  >{\raggedright\arraybackslash}p{(\linewidth - 4\tabcolsep) * \real{0.3333}}
  >{\raggedright\arraybackslash}p{(\linewidth - 4\tabcolsep) * \real{0.3333}}@{}}
\toprule\noalign{}
\begin{minipage}[b]{\linewidth}\raggedright
Closure
\end{minipage} & \begin{minipage}[b]{\linewidth}\raggedright
Group structure
\end{minipage} & \begin{minipage}[b]{\linewidth}\raggedright
Consistency
\end{minipage} \\
\midrule\noalign{}
\endhead
\bottomrule\noalign{}
\endlastfoot
G6 (PAPER\_1159) & \(\Phi_{\rm res} = 5/6 = (D-1)/D\) at
\(D_{\rm BSFG}=6\) & \(D_{\rm BSFG}=6\) from \(26\to 10\to 6\) chain \\
G7 (PAPER\_1160) & \(F_{\rm TRZ} = 1/10 = 1/|SO(5)|\) &
\(SO(5)\subset SO(6)\subset SO(26)\) \\
G8 (PAPER\_1161) & \(26!\) Pochhammer at \(D_{\rm crit}=26\) &
\(D_{\rm crit}=26\) root of all branches \\
\textbf{G3 (this paper)} & \(SO(2) \subset SO(26)\) light-cone embedding
& \(D_{\rm crit}=26\) root \\
\end{longtable}
}

All four closures (G3, G6, G7, G8) descend from the same single integer
\(D_{\rm crit} = 26\) -- the bosonic critical dimension. The embedding
\(SO(26) \supset SO(24) \times SO(2)\) is the \textbf{organising group}
under which: * SO(2) supplies the DPM (G3); * SO(24) contains SO(5),
SO(6), SO(2,3) used in F\_TRZ, \(\Phi_{\rm res}\) closures (G6, G7); *
The 26 of SO(26) carries the 26-fold Pochhammer derivative (G8).

\subsection{8. Result -- G3 Closure
Statement}\label{result-g3-closure-statement}

\textbf{Theorem (G3 Closure).} Let
\(\mathcal{L}_{\rm DPM} = -\tfrac{1}{4}F^{\rm DPM}_{[12]}F^{[12]\rm DPM}\)
denote the DPM Yang-Mills Lagrangian with \(F^{\rm DPM}\) the
field-strength of \(SO(2)\) generated by \(T_{[12]}\). Let
\(\iota: SO(2)
\hookrightarrow SO(26)\) be the embedding via the first two coordinates
\((X^1, X^2)\). Then:

\begin{enumerate}
\def\labelenumi{\arabic{enumi}.}
\tightlist
\item
  \(\iota\) coincides with the light-cone gauge embedding used in
  standard bosonic string quantization (PAPER\_050).
\item
  \(\iota\) has Dynkin index \(T(\iota) = 1\) (minimal, irreducible).
\item
  The vector \(\mathbf{26}\) branches as
  \((\mathbf{24},\mathbf{1}_0)\oplus(\mathbf{1},\mathbf{1}_{+1})
  \oplus(\mathbf{1},\mathbf{1}_{-1})\), identifying the CW and CCW
  monopoles as the \(q=\pm 1\) SO(2)-charged singlets.
\item
  \(F_{\rm DPM} = I\cdot A\cdot (\omega_1-\omega_2)\) is the SO(2)
  charge-difference current, \(\omega_1 - \omega_2\) being the
  eigenvalue separation of the \(\pm i\) eigenstates of \(T_{[12]}\).
\end{enumerate}

\textbf{Free parameters introduced: ZERO.}

\subsection{9. Updated Catalog Status}\label{updated-catalog-status}

{\def\LTcaptype{none} % do not increment counter
\begin{longtable}[]{@{}
  >{\raggedright\arraybackslash}p{(\linewidth - 4\tabcolsep) * \real{0.3333}}
  >{\raggedright\arraybackslash}p{(\linewidth - 4\tabcolsep) * \real{0.3333}}
  >{\raggedright\arraybackslash}p{(\linewidth - 4\tabcolsep) * \real{0.3333}}@{}}
\toprule\noalign{}
\begin{minipage}[b]{\linewidth}\raggedright
Lagrangian gap
\end{minipage} & \begin{minipage}[b]{\linewidth}\raggedright
Status (pre-250)
\end{minipage} & \begin{minipage}[b]{\linewidth}\raggedright
Status (post-250)
\end{minipage} \\
\midrule\noalign{}
\endhead
\bottomrule\noalign{}
\endlastfoot
G1 (V(UA) polynomial) & OPEN & OPEN \\
G2 (beta\_i index) & OPEN & OPEN \\
\textbf{G3 (DPM SO(2))} & \textbf{OPEN} & \textbf{CLOSED (this paper,
group-theoretic)} \\
G4 (T\^{}22 moduli) & OPEN & OPEN \\
G5 (KK tower) & CLOSED (PAPER\_1162) & CLOSED \\
G6 (Phi\_res ID) & CLOSED (PAPER\_1159) & CLOSED \\
G7 (F\_TRZ ID) & CLOSED (PAPER\_1160) & CLOSED \\
G8 (26! emergence) & CLOSED (PAPER\_1161) & CLOSED \\
\end{longtable}
}

\begin{quote}
\textbf{Session 253 Update (PAPER\_1166):} ALL 8 Lagrangian gaps now
closed. The \(SO(26)\supset SO(24)\times SO(2)\) branching used here
generalises in PAPER\_1164 (G4) to the full \(T^{22}\) moduli
stabilisation, completing the dimensional reduction chain.
\end{quote}

\textbf{Result:} \textbf{5 of 8 gaps closed}; 3 remain (G1, G2, G4). All
five closures (G3, G5, G6, G7, G8) descend from the single integer
\(D_{\rm crit}=26\) via the canonical decomposition
\(SO(26)\supset SO(24)\times SO(2)\).

\subsection{10. Conclusions}\label{conclusions}

\begin{itemize}
\tightlist
\item
  The DPM SO(2) gauge group is the \textbf{light-cone gauge 2-plane} of
  textbook bosonic string light-cone quantization at
  \(D_{\rm crit}=26\).
\item
  The embedding \(\iota: SO(2)\hookrightarrow SO(26)\) has Dynkin index
  1 (minimal, irreducible) and dimensional decomposition
  \(325 = 276 + 1 + 48\) that factorises through \(48 = 24 \times 2\)
  (the textbook GSW transverse \(\times\) light-cone count).
\item
  The CW and CCW monopoles are the \(q=\pm 1\) SO(2)-charged singlets in
  the branching of the vector \(\mathbf{26}\).
\item
  \(F_{\rm DPM} = I\cdot A\cdot(\omega_1-\omega_2)\) is identically the
  SO(2) charge-difference current.
\item
  G3 closed; \textbf{5 of 8 Lagrangian gaps now closed}. Only G1, G2, G4
  remain (V(UA) polynomial coefficients, beta\_i index structure,
  T\^{}22 moduli stabilisation).
\item
  Cumulative impact across PAPERs 1159-1163: zero-mode dominance and
  SO(2) DPM embedding rigorously identified; 6 free numerical inputs
  reduced to \textbf{2 textbook integers} (\(D_{\rm crit}=26\),
  \(D_{\rm BSFG}=6\)).
\end{itemize}

\begin{center}\rule{0.5\linewidth}{0.5pt}\end{center}

\subsection{§SM Anchors (G3 Compliance
Table)}\label{sm-anchors-g3-compliance-table}

{\def\LTcaptype{none} % do not increment counter
\begin{longtable}[]{@{}
  >{\raggedright\arraybackslash}p{(\linewidth - 8\tabcolsep) * \real{0.1875}}
  >{\raggedright\arraybackslash}p{(\linewidth - 8\tabcolsep) * \real{0.1875}}
  >{\raggedleft\arraybackslash}p{(\linewidth - 8\tabcolsep) * \real{0.2500}}
  >{\raggedright\arraybackslash}p{(\linewidth - 8\tabcolsep) * \real{0.1875}}
  >{\raggedright\arraybackslash}p{(\linewidth - 8\tabcolsep) * \real{0.1875}}@{}}
\toprule\noalign{}
\begin{minipage}[b]{\linewidth}\raggedright
Anchor
\end{minipage} & \begin{minipage}[b]{\linewidth}\raggedright
Symbol
\end{minipage} & \begin{minipage}[b]{\linewidth}\raggedleft
Value
\end{minipage} & \begin{minipage}[b]{\linewidth}\raggedright
Source
\end{minipage} & \begin{minipage}[b]{\linewidth}\raggedright
Used in
\end{minipage} \\
\midrule\noalign{}
\endhead
\bottomrule\noalign{}
\endlastfoot
Bosonic critical dim & \(D_{\rm crit}\) & 26 & Polyakov (textbook) &
\(SO(26)\) \\
Transverse Lorentz & \(SO(24)\) & \(\dim = 276\) & GSW Vol.1 §2 &
physical modes \\
DPM gauge & \(SO(2)_{\rm DPM}\) & \(\dim = 1\) & this paper \(\equiv\)
light-cone & CW/CCW pair \\
Coset count & \(SO(26)/(SO(24)\times SO(2))\) & \(48 = 24\times 2\) &
this paper & off-diagonal Lorentz \\
Dynkin index & \(T(\iota)\) & 1 & this paper & minimal embedding \\
Branching of \(\mathbf{26}\) &
\(\to (\mathbf{24},\mathbf{1}_0)\oplus 2\cdot(\mathbf{1},\mathbf{1}_{\pm1})\)
& -- & this paper & CW + CCW + 24 transverse \\
DPM current & \(F_{\rm DPM} = I A(\omega_1-\omega_2)\) & -- &
\href{COMPLETE_UQFF_EQUATIONS_REFERENCE.md\#L549}{COMPLETE\_UQFF\_EQUATIONS\_REFERENCE.md
L549} & SO(2) charge difference \\
\end{longtable}
}

\begin{center}\rule{0.5\linewidth}{0.5pt}\end{center}

\subsection{References}\label{references}

\begin{enumerate}
\def\labelenumi{\arabic{enumi}.}
\tightlist
\item
  Murphy, D. T., ``PAPER\_050: 26D Manifold Compactification to 3+1
  Spacetime'' -- light-cone gauge 26 = 24 + 2.
\item
  Murphy, D. T., ``PAPER\_147: F\_DPM Vortical Resonance'' --
  \(F_{\rm DPM} = IA(\omega_1-\omega_2)\) derivation.
\item
  Murphy, D. T., ``PAPER\_1143: Nambu-Goto Bosonic String SCm 26D'' --
  light-cone normal-ordering.
\item
  Murphy, D. T., ``PAPER\_1159, 1160, 1161, 1162'' -- prior G6, G7, G8,
  G5 closures via \(D_{\rm crit}=26\) flow.
\item
  M. B. Green, J. H. Schwarz, E. Witten, \emph{Superstring Theory} Vol.
  1 (Cambridge UP, 1987) Ch. 2 -- light-cone quantization, transverse
  \(SO(24)\), eq. 2.3.45.
\item
  J. Polchinski, \emph{String Theory} Vol. 1 (Cambridge UP, 1998) §1.3
  -- light-cone gauge fix.
\item
  \href{_lagrangian_rederivation_outline.py\#L56}{\_lagrangian\_rederivation\_outline.py
  L56-58} -- G3 gap statement.
\item
  \href{COMPLETE_UQFF_EQUATIONS_REFERENCE.md\#L549}{COMPLETE\_UQFF\_EQUATIONS\_REFERENCE.md
  L549} -- \(F_{\rm DPM}\) canonical form.
\end{enumerate}

\begin{center}\rule{0.5\linewidth}{0.5pt}\end{center}

\textbf{Signed:} Daniel T. Murphy \textbf{Date:} May 10, 2026 (Session
250) \textbf{Repository:} Star-Magic, commit pending
\textbf{Verification:}
\(\dim SO(26) - \dim SO(24) - \dim SO(2) = 325 - 276 - 1 = 48 = 24\times 2\)
(exact, integer); branching rule of \(\mathbf{26}\) under
\(SO(24)\times SO(2)\) standard textbook result.

\end{document}
